\chapter{Conclusion}\label{chapter:conclusion}
Short introduction. It should relect on the primary and secondary research questions.

\section{Highlights}\label{sec:con-highlights}
Mention the strengths (such as adaptiveness through online learning) and weaknesses (such as more CPU time required) of the framework.

Also elaborate on the main contributions of the thesis.

\section{Limitations}\label{sec:con-limitations}

% --- 1. Hardware ----------------------------------------------------------
Table~\ref{tab:solver-summary} compares runs that differ in hardware, parallelism, and total budget,
and none of its gaps correct for these differences. \citet{queirogaXLInstancesCapacitated2026} ran
every published solver on a single thread of a machine with two AMD EPYC~9654
processors (Zen~4, 2.40\,GHz), 60~seeds of two hours per instance. \gls{drispi} ran
three seeds of two hours per instance on eight pinned cores of two Intel Xeon
Platinum 8280L processors (Cascade Lake, 2.70\,GHz; Section~\ref{sec:stu-implementation}). Both machines were
shared: the published runs executed as up to 50 parallel jobs on one machine,
\gls{drispi} as six concurrent containers, three per socket, pinned to physical cores.

\citet{accorsiGuidelinesComputationalTesting2022} recommend bringing running times from different processors
to a common scale with the PassMark single-thread rating before comparing them.
The EPYC~9654 rates 2,898\footnote{\url{https://www.cpubenchmark.net/cpu.php?cpu=AMD+EPYC+9654&id=5088}, accessed September 19, 2026.} and the Platinum 8280L 2,316\footnote{\url{https://www.cpubenchmark.net/cpu.php?cpu=Intel+Xeon+Platinum+8280L+@+2.70GHz&id=6729}, accessed September 19, 2026.}, a factor of 1.25. That factor is a lower bound: the 8280L rating rests on few samples and is measured at the 4.0\,GHz single-core turbo, which the chip cannot hold with 24 of its 28 cores busy. Scaled this way, \gls{drispi}'s 7,200 seconds
per run correspond to at most 5,800 seconds on the published solvers' processor,
per core. Table~\ref{tab:solver-summary} reports unscaled gaps: \gls{drispi} was not rerun with a longer
limit to offset the slower hardware, and a single-thread rating says nothing
about how well eight cores are used in parallel.

% --- 2. Budget: mean vs best ------------------------------------------------
The budget asymmetry points in opposite directions for the two columns of
Table~\ref{tab:solver-summary}. Per run, \gls{drispi} used 16 core-hours (at most 12.8 after PassMark scaling)
against two CPU-hours for each published run, so the mean-of-$N$ column compares runs in which \gls{drispi} had eight
times the cores, on older hardware. Those eight cores do not deliver eight times
the search. Subcluster routing runs in waves that leave \texttt{dri} cores idle
whenever the last wave holds fewer than six subclusters, the main loop blocks
while it waits for BG-AILS (Figure~\ref{fig:pipeline-swimlane}), and three containers share the memory
bandwidth and last-level cache of each socket.
% TODO: quantify if you have it (e.g. mean dri core utilization per iteration).
Per instance, the relation reverses. A published best-of-60 draws on 120 CPU-hours
and 120 hours of wall-clock time, \gls{drispi}'s best-of-3 on 48 core-hours (at most 38
after PassMark scaling) and six hours. The minimum of 60 runs from a distribution lies systematically below the
minimum of three, so the best-of-$N$ column favors the published solvers by
construction. Neither column is an equal-resource ranking.

% --- 3. Run length and HAOS convergence ------------------------------------
The 7,200-second limit matches the published per-run budget, not the run length
\gls{drispi} needs. HAOS learns online. For the first ten iterations it draws uniformly
and updates nothing (Section~3.1), and after that every draw of a weak arm still
costs a full decomposition, routing, and BG-AILS pass. Under a fixed wall-clock
limit, the number of completed iterations therefore bounds how far the wheels can
move. That bound is tightest where one routing wave is expensive: large $n$ and,
separately, long routes. Section~\ref{sec:stu-convergence} reports the realized iteration counts and
[TODO: share of runs still improving in the final 20 minutes]. Whether a longer
run would close the remaining gap on these instances was not tested.

% --- 4. Statistical resolution (S = 3) ---------------------------------------
$S = 3$ was sized for the grand mean $\bar{G}$ (Section~\ref{sec:stu-repetitions}) and supports claims
at that level only. The forecast 95\% half-width of $\bar{G}$ is 0.013\,pp. For a
single instance, the mean of three seeds at the pilot standard deviation of
0.118\,pp has a standard error of 0.068\,pp and, with two degrees of freedom
($t_{0.975,2} = 4.30$), a 95\% half-width of about 0.29\,pp. That is more than
twenty times wider, and wider still on the high-variance instances. The
per-instance results in Section~\ref{sec:stu-results} (the largest gaps in Table~\ref{tab:gap-tail}, the instances
where \gls{drispi}'s best-of-3 matches or beats a published result in Table~\ref{tab:matches-beats}, and
[TODO: small subgroups in Section~\ref{sec:stu-characteristics}]) are therefore descriptive. Across
100 instances, some differences of that size will arise from the seed alone.

The seed term is also not the largest source of error. Two pilot campaigns that
differed only in execution environment showed a batch effect of about 0.039\,pp,
three times the $S = 3$ half-width, and more seeds would not have reduced it.
Shared-server load and Docker pinning are part of the measurement rather than a
nuisance the design controlled.

% --- 5. No fixed-configuration control arm ----------------------------------
The contribution of HAOS is established descriptively, through the weight
trajectories of Section~\ref{sec:stu-haos}, and not against a control arm that freezes the draw
at one configuration for the whole run under the same budget. That arm was not
run. The 300-run campaign behind Section~\ref{sec:stu-comparison} used the server allocation
available for this thesis, and a controlled comparison needs a full campaign run
alongside it; a subset carved out of the existing allocation would confound the
two. \citet{altendeiteringDRSCICVRP2026} leave the same gap for DRSCI, a \gls{drispi}
predecessor that draws paradigm, $k$, and solver uniformly at random each
iteration without learning. Whether HAOS pays a sampling cost relative to a single
committed configuration, or improves on it, remains open.

\section{Future Research}\label{sec:con-future}
Use the HAOS weights to train an offline learning model and inject the results back into the framework. This effectively eliminates the warm-up period and will lead to a substential speed-up of the framework.

Also mention that the framework, other than monolithic solvers, is highly adaptable to different instances and VRP variants. Since it only requires the plug-and-play of capable clustering methods and solvers on the their respective stages.